% TOP_PROBLEM: {"id": 2514, "title": "Kourovka Notebook Problem 21.5", "category_id": 17, "category": {"id": 17, "name": "group_theory", "display_name": "Group Theory", "description": "Problems about groups, group actions, representations, and related algebraic structures.", "slug": "group-theory", "order_index": 17, "created_at": "2026-07-31T15:26:25.671Z"}, "status": "open", "problem_number": "KOU-21.5", "set_id": 10, "statusQualification": "Parts (a) and (b) are marked solved in primary v47; the local-conjugacy part (c) is retained."}
% FIRST_POSTED: 2026-10-01
% ENTRY_KIND: statement_only
% UnsolvedMath ID 2514; source catalogue code KOU-21.5
% !TeX program = lualatex
% Definition and sources only; this file is not an LLM solution attempt.
\documentclass[11pt,a4paper]{article}
\usepackage[margin=25mm]{geometry}
\usepackage{fontspec}
\setmainfont{lmroman10-regular.otf}[BoldFont=lmroman10-bold.otf,
 ItalicFont=lmroman10-italic.otf,BoldItalicFont=lmroman10-bolditalic.otf]
\usepackage{amsmath,amssymb,mathtools}
\usepackage{unicode-math}
\setmathfont{latinmodern-math.otf}
\AtBeginDocument{\let\setminus\smallsetminus}
\usepackage{xurl,hyperref}
\hypersetup{colorlinks=true,urlcolor=blue}
\setcounter{secnumdepth}{0}
\setlength{\parindent}{0pt}
\setlength{\parskip}{5pt}
\setlength{\emergencystretch}{3em}
\begin{document}
\section*{Kourovka Notebook Problem 21.5}\label{2514}
{\small UnsolvedMath ID 2514 · KOU-21.5 · Group Theory · Kourovka Notebook - New Problems, Issue 21}
\subsection{Definitions and mathematical statement}
Let $\Omega$ be a nonempty set. The finitary symmetric group
$\operatorname{FSym}(\Omega)$ consists of permutations moving only finitely
many points. Every finitely generated subgroup of it is finite.
Fix a prime $p$, a transitive subgroup $G\le\operatorname{FSym}(\Omega)$,
a normal subgroup $N\trianglelefteq G$, and a transitive Sylow
$p$-subgroup $S$ of $G$.

Here transitivity means that any point can be sent to any other point.
For these locally finite groups a $p$-subgroup is a subgroup all of whose
elements have order a power of $p$, and a Sylow $p$-subgroup means one
maximal under inclusion. It need not be finite if $\Omega$ is infinite.
The questions are:
\begin{enumerate}
\item Is $S\cap N$ a Sylow $p$-subgroup of $N$?
\item Is $SN/N$, the image of $S$ in $G/N$, a Sylow $p$-subgroup of $G/N$?
\item Are any two transitive Sylow $p$-subgroups of $G$ locally conjugate?
\end{enumerate}
For the last condition, $X,Y\le G$ are locally conjugate if an
automorphism $\alpha\in\operatorname{Aut}(G)$ maps $X$ onto $Y$ and,
for every finite subset $F\subseteq G$, some $g_F\in G$ satisfies
$\alpha(a)=g_F^{-1}ag_F$ for all $a\in F$.

The conjugator may change with $F$, but the automorphism $\alpha$ must
be one fixed map on all of $G$. Thus local conjugacy does not require
one global inner automorphism. All three questions retain the special
transitivity assumption on the Sylow subgroups.

The primary edition records affirmative solutions of parts (a) and (b).
They are retained to define the original multipart entry; part (c) is
the remaining question.
\subsection{Short English statement}
Do transitive Sylow subgroups of finitary permutation groups behave well under normal intersections, quotient images, and local conjugacy?
\subsection{Sources}
\begin{itemize}
\item[S1] ULAM.AI and the UnsolvedMath Contributors, supplied problems.json, record 2514 (KOU-21.5), Kourovka Notebook - New Problems, Issue 21. Catalogue identity and source transcription; CC BY 4.0.
\url{https://huggingface.co/datasets/ulamai/UnsolvedMath}.
\item[S2] The Kourovka Notebook, 21st edition, new problems, Problem 21.5. Expanded formulation of the supplied catalogue statement.
\url{https://arxiv.org/pdf/1401.0300v47}.
\end{itemize}
\end{document}
